\section{Limitations}

\begin{itemize}
  \item We made every score-based choice on the 2025 pilots, as the rules require.
    They have 28 to 59 systems each, so the 2026 results may differ from what the
    pilots suggest.
  \item The organizers re-scored the pilots after we had made our choices. The
    ensemble for rag26 looks better under the earlier scoring than under the
    current one.
  \item On dragun-repgen our local runs used different topic text from the copy on
    TIRA. We do not know whether the test sets on TIRA match the released files,
    and our data submissions used the released files.
  \item \judge{NuggetRecall} depends on the LLM's own knowledge to write the answer
    key, and on the request being a real information need. It never sees
    citations, so it cannot judge attribution.
  \item \judge{Corroborate} cannot reward a correct fact that no other team stated.
    Consensus is a proxy for importance, not a measure of it, and a correct but rare
    fact gets little weight. We did not run the ablations that would show how much
    the leave-one-team-out weight and per-team voting each change the ranking.
  \item Without truth on the test sets we could check only that our leaderboards
    were complete and that every measure varied across systems. We ran no length
    diagnostics on them. A script may compute such a statistic, but the policy
    treats the result as restricted, and it could have steered choices that must
    come from the pilots.
\end{itemize}

\section{Data Handling and Reproducibility}
\label{sec:disclosure}

Development followed the AutoJudge data policy. On the test sets we read response
content only for the first ten topics of each topics file, in file order, and every
choice we made by score came from the pilots. Every bug but one reproduced on a
pilot or on the synthetic dataset that ships with the starter kit, so we never
needed the policy's scrubbing tool. The exception is the exhausted reply budget of
Section~\ref{sec:failures}. It also occurs on the pilots, at lower rates, but we
noticed it on rag26 and saw it there only as counts: unreadable judgments in our
logs, and requests that vLLM stopped at the length limit. We read no response text
to diagnose it. The answer keys of \judge{NuggetRecall} are built from the topic
alone, so they may be read for any topic. The claim banks of \judge{Corroborate}
are built from responses, and we inspected them only on the pilots.

We also submitted two generation runs to TREC 2026 RAG, so two of the 83 rag26 runs
are probably ours. We made no attempt to identify them, by comparison with our own
submissions or by any other means. Neither judge knows which runs are ours, and
both treat every run alike. \judge{NuggetRecall} judges each response on its own
against a key built from the topic. \judge{Corroborate} drops a team's own vote
from each claim's weight, so a team's runs cannot raise the weight of the claims
they are scored on.

Both judges are available as code submissions on TIRA. The source code, workflows
and tests are at \pending{public repository URL}. \judge{NuggetRecall} reproduces
from its prompt cache with the endpoint switched off. \judge{Corroborate} makes no
network calls, and on the same inputs and hardware it produces byte-identical
output.

\section{Summary}

We built two judges for TREC 2026 AutoJudge that fail in different places.
\judge{NuggetRecall} writes an answer key from the question and does best on nugget
coverage. \judge{Corroborate} trusts what many teams agree on, does best on F1 for
long reports and fails on short answers. The largest effect we measured was the
choice of judge model. The most useful lessons, though, concern failure. An LLM
judge can lose most of its judgments without an error and still produce a
leaderboard that looks plausible. A check that measures vary across systems, retry
counts watched on the real data, and a liveness check beside every reproducibility
check caught most such failures. An audit aimed at exactly this kind of failure
caught the rest.
